% ---------------------------------------------------------------------------
% Results fragment for Appendix D, closing subsection "Reconciliation with the earlier
% release" (subsec:corr_reconciliation). It follows Table~\ref{tab:corrections_reconciliation}
% and the "After the rerun, ..." paragraph (both filled from this theme's tbd fills) and adds no
% float with a label that already exists in the sections. New labels:
%   fig:reconcile_scatter   -> figures/reconcile_scatter.pdf
%   fig:reconcile_cells     -> figures/reconcile_cells.pdf
%   fig:reconcile_movers    -> figures/reconcile_movers.pdf
%   fig:reconcile_si_terms  -> figures/reconcile_si_terms.pdf
%   tab:reconcile_coverage  -> tables/reconcile_coverage.tex (\input below)
%   tab:reconcile_movers    -> tables/reconcile_movers.tex   (\input below)
%   tab:reconcile_allied    -> tables/reconcile_allied.tex   (\input below)
%   tab:reconcile_network   -> tables/reconcile_network.tex  (\input below)
% Numbers: paper/results/reconcile.json, written by scripts/paper/fig_reconcile.R from the
% shipped release of 3 September 2026 (main checkout, read only) and the final rerun of
% 12 October 2026 (network at lambda = 0.78). "Shipped" is the September 2026 release, "final"
% the rerun. Rank statistics are over cells rankable in the final release unless stated.
% ---------------------------------------------------------------------------

\paragraph{Matching and coverage.} The comparison uses the headline tables of the two releases,
which cover the eleven technologies of the paper. Every country name of the shipped release is
mapped to its canonical name with the corrected standardiser (item C1), units without an ISO3
code are dropped (item C2), and where two spellings of one economy give the same
country--technology--stage key the record under the canonical spelling is kept. Pillar cells are
matched on (country, technology, stage) and Strategic Index cells on (country, sector).
Table~\ref{tab:reconcile_coverage} reports the counts. Harmonisation alone reduces the shipped
Energy Security table from 270 country strings to 237 economies: it drops 17 rows of two units
without a code (Northern Cyprus and the USSR) and merges 162 rows of 29 economies that appeared
under two or three spellings, such as Viet Nam and Vietnam or Slovakia and the Slovak Republic.
The shipped Economic Opportunity table loses 56 rows of three such units, among them ``Global'',
and merges 234 rows. The final release scores 219 economies in each pillar, each under one name.

Of the harmonised $\ES$ cells, 5{,}522 are in both releases, 45 only in the shipped one and 710
only in the final one. The 45 are the Solar and Wind downstream cells of 18 small territories
outside the World Bank country list, such as the Falkland Islands, Jersey and Vatican City (35
cells), and the Geothermal upstream and midstream cells of five economies, pairs that are no
longer scored in $\ES$ (10 cells). The 710 are new coverage, not relabelled cells: 109 fill Green
Hydrogen midstream, which the shipped $\ES$ did not score, and 601 add economies to existing
cells, most of them midstream cells reached by the rebuilt trade themes (items T1--T5); 432 of
the 710 are rankable. For $\EO$, 6{,}598 cells are in both releases, the same 35 territory cells
are lost, and 533 cells are added, all of them economies newly scored in cells the shipped
release already scored. These additions rest mostly on imputed values (the median imputed weight
share is 0.68), and only 11 of them are rankable. They enlarge the population of each cell more
than they add information about the economies they bring in.

\begin{table}[htbp]
  \centering
  \begin{threeparttable}
  \caption{Coverage of the headline tables: shipped and final release}
  \label{tab:reconcile_coverage}
  \small
  \begin{tabular}{@{}l r r@{}}
    \toprule
    & Shipped release & Final release \\
    \midrule
    \multicolumn{3}{@{}l}{\emph{Energy Security}} \\
    Rows (country--technology--stage cells) & 5{,}746 & 6{,}232 \\
    \quad rankable & -- & 4{,}163 \\
    Distinct country strings & 270 & 219 \\
    \quad non-country units dropped (rows) & 17 & -- \\
    \quad merged spellings (rows) & 162 & -- \\
    Economies after harmonisation & 237 & 219 \\
    Technology--stage pairs scored & 31 & 30 \\
    Cells in both releases & 5{,}522 & 5{,}522 \\
    \quad rankable in the final release & -- & 3{,}731 \\
    Cells in one release only & 45 & 710 \\
    \addlinespace
    \multicolumn{3}{@{}l}{\emph{Economic Opportunity}} \\
    Rows (country--technology--stage cells) & 6{,}923 & 7{,}131 \\
    \quad rankable & -- & 4{,}301 \\
    Distinct country strings & 272 & 219 \\
    \quad non-country units dropped (rows) & 56 & -- \\
    \quad merged spellings (rows) & 234 & -- \\
    Economies after harmonisation & 237 & 219 \\
    Technology--stage pairs scored & 33 & 33 \\
    Cells in both releases & 6{,}598 & 6{,}598 \\
    \quad rankable in the final release & -- & 4{,}290 \\
    Cells in one release only & 35 & 533 \\
    \addlinespace
    \multicolumn{3}{@{}l}{\emph{Strategic Index}} \\
    Scored cells & 571 & 3{,}154 \\
    Countries & 18 & 174 \\
    Sectors & 32 & 30 \\
    Cells in both releases & 477 & 477 \\
    \bottomrule
  \end{tabular}
  \begin{tablenotes}[flushleft]\footnotesize
    \item \emph{Notes:} Headline tables of the eleven technologies. Shipped release: \texttt{index\_outputs.rds} of 3~September 2026; final release: the same file of the final 12~October 2026 run. Shipped country names are mapped to canonical names with \texttt{standardize\_country\_names()} (item C1); units without an ISO3 code (``Global'', ``USSR'', ``Northern Cyprus'') are dropped (item C2) and, where two spellings give the same cell, the record under the canonical spelling is kept. ``Rankable'': the economy is sovereign and at most half of the cell's index weight is imputed (\texttt{config/rankability.yml}); the shipped release had no such rule. ``Cells in one release only'' counts shipped cells absent from the final release (first column) and final cells absent from the shipped release (second column). A Strategic Index cell of the final release requires a Policy index (IEA PAMS, combined with CAT where CAT rates the country) and a rankable EO and ES value; a shipped cell required the four-component Policy index of February 2026 and imputed missing pillar values (item P7).
  \end{tablenotes}
  \end{threeparttable}
\end{table}


Three rules apply to every pillar statistic below. First, a technology--stage cell in which a
pillar is constant across countries in either release carries no information on the ordering of
countries and is excluded. This removes $\ES$ in Nuclear upstream, which was zero for all 217
economies in the shipped release, whose only category was a manufactured zero Investment cell
(item E7), and the shipped Geothermal upstream and midstream cells, also constant at zero. In the
final release Nuclear upstream varies, being scored on uranium resources and production from the
Red Book of the OECD Nuclear Energy Agency and the International Atomic Energy Agency, Trade and
Energy Prices. No $\EO$ cell is constant in either
release. Second, rank statistics are computed over the cells that are rankable in the final
release (Section~\ref{subsec:norm_diagnostics}): 3{,}537 of the 5{,}305 matched non-constant
$\ES$ cells and 4{,}290 of the 6{,}598 $\EO$ cells; the shipped release had no such rule, and the
statistics over all matched cells are reported alongside. Third, correlations are Spearman rank
correlations unless they are labelled Pearson.

\paragraph{Agreement of the pillars.} Figure~\ref{fig:reconcile_scatter} plots the final value
of each cell against the shipped one. Over the 3{,}537 rankable $\ES$ cells of 194 economies the
Spearman correlation is 0.81 (Pearson 0.81); over the 4{,}290 rankable $\EO$ cells of 178
economies it is 0.67 (Pearson 0.69). Over all matched cells the correlations are 0.84 and 0.71.
Agreement differs by stage, and in opposite directions for the two pillars. For $\ES$ it is
highest upstream (0.83) and downstream (0.80) and lowest at the midstream stage (0.55); for
$\EO$ it is lowest upstream (0.50) and highest downstream (0.88), with the midstream stage at
0.66. The decomposition reported below locates the source of the weak stages: in the $\ES$
midstream the reordering comes mainly from the Trade category, which items T1--T5 rebuilt from
all-reporter Comtrade dyads and the reviewed baskets; in the $\EO$ upstream it comes from
categories that enter the cell, among them the Trade category of the new upstream baskets, and
from the rebuilt Trade and Investment inputs.

Levels moved as well as ranks. The mean change, final minus shipped, is $+0.021$ in $\ES$ and
$+0.012$ in $\EO$, but $+0.081$ in the $\ES$ upstream. Part of that shift comes from the 89
upstream cells of 21 economies that were exactly zero in the shipped release, the vertical stripe
at the left edge of the first panel: the shipped aggregation counted an unobserved category as
zero at full weight (item A1), so a cell whose categories were all unobserved or zero scored zero;
the same cells average 0.20 in the final release. The other vertical stripes are groups of
economies that shared one shipped value because they shared imputed category values. A change of
more than 0.1 in absolute value affects 31 percent of the rankable $\ES$ cells (40 percent at the
midstream stage) and 17 percent of the $\EO$ cells (22 percent upstream, 5 percent downstream);
a change of more than 0.2 affects 6.4 and 0.3 percent.

\begin{figure}[htbp]
  \centering
  \includegraphics[width=\textwidth]{reconcile_scatter}
  \caption{Shipped release against final release, cell by cell. Each point is a
  country--technology--stage cell (Energy Security, Economic Opportunity) or a country--sector
  cell (Strategic Index) present in both releases, with the shipped value (3~September 2026) on
  the horizontal axis and the final value on the vertical axis; points above the 45-degree line
  gained. Colour marks the supply-chain stage of cells that are rankable in the final release
  (sovereign economy, at most half of the index weight imputed); grey points are matched cells
  that are not rankable. Energy Security: 5{,}305 matched cells of 219 economies, 3{,}537 of
  them rankable; Economic Opportunity: 6{,}598 and 4{,}290; Strategic Index: the 477 cells of the
  18 members of the shipped hand-picked list that are scored in both releases. Energy Security is
  oriented so that higher values are more secure in both releases. Cells of technology--stage
  pairs in which a pillar is constant across countries in either release (Energy Security in
  Nuclear upstream, and in the shipped Geothermal upstream and midstream) are excluded. Spearman
  correlations over rankable cells: Energy Security 0.81, Economic Opportunity 0.67, Strategic
  Index 0.74. Shipped country names are mapped to canonical names before matching (item C1).
  Sources: \texttt{index\_outputs.rds} of the two releases; \texttt{scripts/paper/fig\_reconcile.R}.}
  \label{fig:reconcile_scatter}
\end{figure}

\paragraph{Level shifts and reorderings.} A pooled correlation mixes shifts in the level of a
whole technology--stage cell with reorderings of countries within it.
Figure~\ref{fig:reconcile_cells} separates the two: for each cell it plots the mean change
against the Spearman correlation, across the rankable economies present in both releases, of the
shipped and final values. The median within-cell correlation is 0.83 for $\ES$ (interquartile
range 0.74--0.91) and 0.77 for $\EO$ (0.61--0.88); 11 of the 28 $\ES$ cells and 20 of the 33
$\EO$ cells fall below 0.8, and four $\EO$ cells below 0.5. Measured by positions rather than
correlations, 45 percent of $\ES$ values and 56 percent of $\EO$ values move by more than ten
places in the ranking of their cell, and 41 and 48 percent by more than a tenth of the economies
ranked in it; the median move is 9 places for $\ES$ and 13 for $\EO$, in cells of 68 to 191 and 40
to 175 rankable economies. The ordering of economies within a cell is therefore not preserved
closely enough for a shipped ranking to stand in for a final one.

To attribute the changes, write a cell's pillar value as $\sum_k \omega_k c_k$, with $c_k$ the
score of category $k$ and $\omega_k$ its normalised weight, an unobserved shipped category
counting as $c_k=0$ at full weight, which is how the shipped value was computed (item A1). A
two-factor Shapley decomposition splits the change into a \emph{score effect} for each category
present in both releases,
$\tfrac12(\omega_k^{\text{shipped}}+\omega_k^{\text{final}})(c_k^{\text{final}}-c_k^{\text{shipped}})$,
and a \emph{composition effect}, the remainder, which comes from categories that enter or leave
the cell and from the renormalisation of the weights. The decomposition reconstructs both
releases' values exactly in every matched cell. Composition effects account for 25 percent of the
summed absolute change in $\ES$ and 40 percent in $\EO$. Among score effects, Consumption (21
percent), Trade (19 percent) and Foreign Dependency (16 percent) dominate in $\ES$, and
Technological Readiness (18 percent), Cost Competitiveness (12 percent), Technology Demand and
Trade (9 percent each) in $\EO$. The two kinds of effect play different roles, and the difference
follows from the construction. When the category that enters or leaves is constant within a
cell, such as the zero Investment cells, Energy Prices or Technological Readiness, its entry or
exit applies one increasing affine transformation to the values of the countries that share the
cell's category set, so it moves the level of the cell and leaves the ordering within it
unchanged. The composition effect is the largest mean effect in 6 of the 9 $\ES$ upstream cells
and 7 of the 11 $\EO$ upstream cells, and Technological Readiness, now mapped by stage (item T7),
sets the level in 7 of the 11 $\EO$ downstream cells. The category whose effect varies most across
countries, which is what reorders them, is instead a rebuilt input: Foreign Dependency in 6 of the
9 $\ES$ upstream cells, Trade in 6 of the 9 $\ES$ midstream cells and in 4 of the 11 $\EO$
upstream and midstream cells each, and Consumption or Investment in 9 of the 10 $\ES$ downstream
cells.

At the bottom of Figure~\ref{fig:reconcile_cells} lie cells in which the ranking rests heavily on
imputed values in at least one release. The extreme case is $\EO$ in Green Hydrogen midstream,
where the within-cell correlation over its 40 rankable economies is $-0.55$. In the shipped
release only Cost Competitiveness varied across economies in this cell, and it was imputed in 55
percent of its 110 records. In the final release Foreign Dependency enters, measured by
electrolyser capacity (item E8), but it is imputed for 86 percent of the 213 economies scored and
Cost Competitiveness for 84 percent. Neither ranking in this cell should be read as a measurement
of the economies' positions. The next lowest cells are $\EO$ in Green Hydrogen upstream (0.34),
Electric Grid upstream (0.36), where the review replaced electricity with a copper and aluminium
basket, and Solar upstream (0.39), which gained the Trade category of the silicon-metal basket.
In $\ES$ the lowest is Electric Grid upstream (0.58). $\ES$ in Nuclear midstream combines the
largest level shift, a mean gain of 0.24, with a within-cell correlation of 0.59: the shift comes
from the removal of the zero Investment category, so that the cell's $\ES$ is now its Trade score,
and the reordering from the rebuilt Trade category itself.

\begin{figure}[htbp]
  \centering
  \includegraphics[width=\textwidth]{reconcile_cells}
  \caption{Level shift and rank agreement by technology--stage cell. Each point is one
  technology--stage cell. Horizontal axis: mean change of the pillar across the economies that
  are scored in the cell in both releases and rankable in the final release (final minus shipped,
  unit-interval scale). Vertical axis: Spearman correlation, across the same economies, of the
  shipped and final values. Colour marks the stage. Labels name the three cells with the lowest
  correlation and the three with the largest absolute mean change in each panel. Energy Security:
  28 cells (Nuclear upstream was constant in the shipped release and is excluded; Geothermal
  upstream and midstream and Electric Grid downstream have no final score; Green Hydrogen
  midstream had no shipped one); Economic Opportunity: all 33 cells. A point far from zero on the
  horizontal axis but high on the vertical axis is a level shift that left the ordering of
  economies largely intact. Sources: as in Figure~\ref{fig:reconcile_scatter}.}
  \label{fig:reconcile_cells}
\end{figure}

\paragraph{Largest changes among coalition members.} Restricted to the 57 members of the final
coalition, the rankable matched cells number 1{,}468 for $\ES$ and 1{,}716 for $\EO$, with
Spearman correlations of 0.82 and 0.68 and changes larger than 0.1 in 28 and 18 percent of cells,
close to the full-sample figures. Figure~\ref{fig:reconcile_movers} shows the twelve largest
absolute changes in each pillar, with at most two per technology--stage cell so that one
cell-wide change does not fill the list, and Table~\ref{tab:reconcile_movers} gives their
decomposition. Uncapped, the twelve largest $\ES$ changes are all nuclear cells, ten of them in
Nuclear upstream, with gains of 0.38 to 0.74. The two largest, Canada and Namibia in Nuclear
upstream, rise from zero to 0.74 and 0.72: the shipped cell held only the zero Investment
category, and the final cell scores the uranium resources and production of the Red Book, which
place both among the leading holders and producers, together with Trade and Energy Prices.
Germany and Croatia in Nuclear midstream gain 0.38 because the zero Investment category, with a
weight of 0.43, has gone. The other movers trace the rebuilt themes: Türkiye in Oil downstream
($+0.32$) and New Zealand in Green Hydrogen downstream ($-0.24$) through the consumption rebuild
(item E5); Türkiye and Czechia in Batteries downstream ($-0.27$ and $-0.25$) through an
Investment value that is now zero-filled (item E7) and the consumption rebuild; Thailand and
Ghana in Wind upstream and Gabon and Oman in Solar upstream (gains of 0.26 to 0.30) through the
removal of the zero Investment category, the entry of Energy Prices and of the Trade category of
the new upstream baskets, and a rebuilt Production or Foreign Dependency score. Gabon did not
report imports in 2024, so its new Trade category rests on its partners' mirror declarations. In
$\EO$ the largest changes are smaller ($-0.20$ to $+0.25$ uncapped): Canada and Namibia in
Nuclear upstream again ($+0.25$ and $+0.24$, from the Red Book production data, the entry of
Energy Prices and Trade and the removal of the zero Investment category); Slovakia in Electric
Vehicles midstream ($+0.23$), whose records were split between two spellings (item C1) and whose
Foreign Dependency and Investment categories are now observed; Germany and Portugal in Green
Hydrogen midstream ($+0.23$ and $+0.20$, item T6 and the new Foreign Dependency category, item
E8); and the United States in Electric Grid upstream ($+0.18$), where the reviewed copper and
aluminium basket rebuilt Trade. Eight of the twelve $\ES$ changes and seven of the twelve $\EO$
changes are driven mainly by the composition effect. The largest changes among the members are
thus mostly level shifts of whole cells, caused by the removal of manufactured zeros and the
entry of new categories and data. They matter most for comparisons across the cells of one
country, which is what the Strategic Index makes, and less for comparisons of countries within a
cell.

\begin{figure}[htbp]
  \centering
  \includegraphics[width=\textwidth]{reconcile_movers}
  \caption{The twelve largest changes among coalition members, by pillar, with at most two per
  technology--stage cell. Each row is one member's technology--stage cell, labelled by ISO3 code.
  Hollow circles mark the shipped value (3~September 2026) and filled circles the final value,
  coloured by stage; the segment spans the change. The text at the right names the attributed
  cause in compact form: the largest effect in a two-factor Shapley decomposition of the change
  into category score effects and a composition effect (categories entering or leaving the cell
  and the renormalisation of the weights), followed by the second largest when it is at least half
  as large, with the data change behind it (``Red Book uranium'': OECD-NEA/IAEA uranium resources
  and production; ``mirror'': trade measured from partners' declarations). Item codes refer to
  Table~\ref{tab:corrections_index}; the full decomposition is in
  Table~\ref{tab:reconcile_movers}. Sample: the 57 members of the final coalition, cells present in
  both releases and rankable in the final release, constant cells excluded (Energy Security in
  Nuclear upstream, constant at zero in the shipped release, is included as a mover). Sources:
  \texttt{category\_contributions} of the two \texttt{index\_outputs.rds}.}
  \label{fig:reconcile_movers}
\end{figure}

\begingroup
  \singlespacing\scriptsize
  \setlength{\tabcolsep}{3pt}
  \setlength{\LTcapwidth}{\textwidth}
  \begin{longtable}{@{}>{\raggedright\arraybackslash}p{0.75in} >{\raggedright\arraybackslash}p{0.95in} c c c >{\raggedright\arraybackslash}p{1.05in} >{\raggedright\arraybackslash}p{1.9in}@{}}
  \caption{Largest pillar changes among coalition members and their attribution}
  \label{tab:reconcile_movers}\\
    \toprule
    Member & Cell & Shipped & Final & \makecell{Composition\\effect} & Largest score effect & Attributed cause \\
    \midrule
  \endfirsthead
    \multicolumn{7}{@{}l}{\small\textbf{Table \thetable:} continued}\\
    \toprule
    Member & Cell & Shipped & Final & \makecell{Composition\\effect} & Largest score effect & Attributed cause \\
    \midrule
  \endhead
    \bottomrule
  \endfoot
    \multicolumn{7}{@{}l}{\emph{Energy Security}} \\
    Canada & Nuclear Upstream & 0.000 & 0.741 & +0.741 & -- & Zero Investment cell removed (E7); uranium Reserves and Production added (Red Book); Energy Prices added (A4); Trade added (T1--T5; all-reporter dyads) \\
    Namibia & Nuclear Upstream & 0.000 & 0.716 & +0.716 & -- & Zero Investment cell removed (E7); uranium Reserves and Production added (Red Book); Energy Prices added (A4); Trade added (T1--T5; all-reporter dyads) \\
    Germany & Nuclear Midstream & 0.345 & 0.729 & +0.286 & Trade +0.099 & Zero Investment cell removed (E7) \\
    Croatia & Nuclear Midstream & 0.288 & 0.665 & +0.251 & Trade +0.126 & Zero Investment cell removed (E7) + Trade rebuilt (T1--T5; all-reporter dyads) \\
    Turkiye & Oil Downstream & 0.468 & 0.787 & +0.000 & Consumption +0.320 & Consumption rebuilt (E5) \\
    Thailand & Wind Upstream & 0.112 & 0.415 & +0.127 & Production +0.123 & Zero Investment cell removed (E7); Energy Prices added (A4); Trade added (T1--T5; all-reporter dyads) + Production rebuilt (E1, E3) \\
    Ghana & Wind Upstream & 0.018 & 0.296 & +0.105 & Minerals Trade +0.088 & Zero Investment cell removed (E7); Energy Prices added (A4); Trade added (T1--T5; all-reporter dyads) + Minerals Trade observed, was imputed (E6) \\
    Turkiye & Batteries Downstream & 0.580 & 0.312 & +0.000 & Investment $-$0.165 & Investment zero-filled (E7) + Consumption rebuilt (E5) \\
    Gabon & Solar Upstream & 0.000 & 0.263 & +0.105 & Foreign Dependency +0.086 & Zero Investment cell removed (E7); Energy Prices added (A4); Trade added (T1--T5; all-reporter dyads, mirror imports) + Foreign Dependency rebuilt (E2--E3) \\
    Oman & Solar Upstream & 0.000 & 0.262 & +0.090 & Foreign Dependency +0.086 & Zero Investment cell removed (E7); Energy Prices added (A4); Trade added (T1--T5; all-reporter dyads) + Foreign Dependency rebuilt (E2--E3) \\
    Czech Republic & Batteries Downstream & 0.580 & 0.331 & +0.000 & Investment $-$0.165 & Investment zero-filled (E7) + Consumption rebuilt (E5) \\
    New Zealand & Green Hydrogen Downstream & 0.451 & 0.212 & +0.000 & Consumption $-$0.238 & Consumption rebuilt (E5) \\
    \addlinespace
    \multicolumn{7}{@{}l}{\emph{Economic Opportunity}} \\
    Canada & Nuclear Upstream & 0.398 & 0.647 & +0.266 & Technology Demand $-$0.015 & Zero Investment cell removed (E7); Energy Prices added (A4); Production added (E1, E3; uranium, Red Book); Trade added (T1--T5; all-reporter dyads) \\
    Namibia & Nuclear Upstream & 0.379 & 0.616 & +0.257 & Cost Competitiveness $-$0.016 & Zero Investment cell removed (E7); Energy Prices added (A4); Production added (E1, E3; uranium, Red Book); Trade added (T1--T5; all-reporter dyads) \\
    Slovak Republic & Electric Vehicles Midstream & 0.358 & 0.591 & +0.000 & Foreign Dependency +0.148 & Spellings merged (C1); Foreign Dependency observed + Investment observed \\
    Germany & Green Hydrogen Midstream & 0.278 & 0.503 & +0.066 & Cost Competitiveness +0.098 & Cost Competitiveness rebuilt (T6) + Foreign Dependency added (E8) \\
    Portugal & Green Hydrogen Midstream & 0.278 & 0.479 & +0.070 & Cost Competitiveness +0.118 & Cost Competitiveness rebuilt (T6) + Foreign Dependency added (E8) \\
    United States & Wind Upstream & 0.397 & 0.584 & +0.164 & Production +0.064 & Zero Investment cell removed (E7); Energy Prices added (A4); Trade added (T1--T5; all-reporter dyads) \\
    United States & Electric Grid Upstream & 0.427 & 0.608 & +0.069 & Trade +0.079 & Trade rebuilt (T1--T5; all-reporter dyads, basket review) + Production rebuilt (E1, E3) \\
    Viet Nam & Coal Downstream & 0.465 & 0.640 & +0.000 & Production +0.052 & Spellings merged (C1); Production observed + Consumption rescored \\
    France & Nuclear Midstream & 0.474 & 0.647 & +0.124 & Trade +0.034 & Zero Investment cell removed (E7) \\
    Colombia & Nuclear Midstream & 0.340 & 0.511 & +0.094 & Trade +0.044 & Zero Investment cell removed (E7) \\
    Norway & Geothermal Midstream & 0.434 & 0.602 & +0.148 & Cost Competitiveness +0.038 & Zero Investment cell removed (E7) \\
    Colombia & Geothermal Midstream & 0.389 & 0.557 & +0.134 & Cost Competitiveness +0.050 & Zero Investment cell removed (E7) \\
  \end{longtable}
  \vspace{-6pt}\footnotesize\noindent
    \emph{Notes:} The 12 largest absolute changes, final minus shipped release, of $\ES_{its}$ and of $\EO_{its}$ among the 57 members of the final coalition, at most 2 per technology--stage cell, over the headline cells present in both releases that the final release ranks (sovereign economy, at most half of the index weight imputed); cells that were constant across countries in the shipped release are included. The change is split by a two-factor Shapley decomposition of the weighted category mean: the score effect of a category present in both releases is $\tfrac12(\omega^{\text{shipped}}_k+\omega^{\text{final}}_k)(c^{\text{final}}_k-c^{\text{shipped}}_k)$, with $\omega_k$ the normalised category weight and an unobserved shipped category counted as $c_k=0$ (item A1); the composition effect is the remainder, due to categories that enter or leave and to the renormalisation of the weights. The cause names the largest effect, and the second when it is at least half as large; item codes refer to Table~\ref{tab:corrections_index}. Data changes: ``uranium, Red Book'' = identified recoverable resources and mine production of the OECD-NEA/IAEA Red Book; ``USGS MCS 2026'' = critical-mineral reserves of the USGS Mineral Commodity Summaries 2026; ``all-reporter dyads'' = trade rebuilt from all-reporter Comtrade dyads (importer report first, else the exporter's declaration $\times 1.06$) and the energy-trade file rebuilt from the same pull; ``basket review'' = a stage whose HS6 basket the October 2026 review changed; ``mirror imports'' = an economy that did not report 2024 imports, whose imports come from partners' declarations; ``TRL by stage'' = technology readiness mapped by stage from the IEA Clean Technology Guide. Source: \texttt{category\_contributions} of the two \texttt{index\_outputs.rds}.\par
\endgroup
\medskip


\paragraph{Strategic Index.} The shipped Strategic Index scored 571 cells of 18 economies of its
hand-picked list in 32 sectors; the final one scores 3{,}154 cells of 174 countries in 30 sectors,
because the Policy index now reaches 215 economies and the rankability rule decides which pillar
cells may enter. On the 477 cells of the 18 economies scored in both releases the Spearman
correlation is 0.74 (Pearson 0.76), the mean absolute change is 0.041 and 4.4 percent of cells
change by more than 0.1. Across sectors within a country the median Spearman correlation is 0.74,
ranging from 0.42 for Argentina to 0.92 for Kenya. Of the 94 shipped cells of these economies
without a final value, 34 lack a final $\ES$ value, most of them in Electric Grid downstream (17)
and Geothermal upstream (16), sectors the shipped release had filled with country means (item
P7), and 60 are excluded because their $\EO$ or $\ES$ cell is not rankable.

The policy term changed most. The shipped Policy index combined four components (IEA PAMS, CAT,
the February 2026 NIPO component and the dual-use score); the final one combines IEA PAMS, which
is required, and CAT where CAT rates the country, and every one of the 18 economies has both.
Between releases the within-country policy term has a Spearman correlation of only 0.61 (Pearson
0.59), and it accounts for 40 percent of the summed absolute change of the index, against 35
percent for the security-need term, 22 percent for $\EO$ and 3 percent for the stage and
technology weights, which differ only through the technology weight of oil (0.12 shipped, 0 final)
and coal (0 and 0.02). Figure~\ref{fig:reconcile_si_terms} splits the twelve largest changes
exactly into the five weighted terms; the policy term is the largest contributor in seven of them
and the security-need term in five. The United States illustrates the policy effect: its policy
term rises from 0.56 to 1 in Batteries downstream and falls from 1 to 0.25 in Oil midstream,
moving the index by $+0.17$ and $-0.15$. The security-need effect is clearest in Nuclear
upstream, where the shipped $\ES$ of zero gave every country the maximum need: with the Red Book
data Australia's index falls by 0.16 and India's by 0.13. Sector means explain 66 percent of the
variance of the index in both releases, and the Spearman correlation within sectors is 0.67, so
the agreement of the Strategic Index owes much to the sector structure that the stage and
technology weights impose on both releases (Section~\ref{sec:strategic}).

\begin{figure}[htbp]
  \centering
  \includegraphics[width=\textwidth]{reconcile_si_terms}
  \caption{The twelve largest changes in the Strategic Index among the economies scored in both
  releases, split into the weighted terms of
  $\SI=0.2\,\POL+0.3\,\EO+0.3\,(1-\ES)+0.1\,\lambda_s+0.1\,g_t$. Each bar is one country--sector
  cell; the segments are the changes of the policy term, the $\EO$ term, the security-need term
  and the two fixed weights (combined), and the diamond is the net change (final minus shipped),
  which the segments sum to exactly. The shipped policy term is the within-country S-curve of the
  four-component Policy index of February 2026 (IEA PAMS, CAT, NIPO and dual use); the final one
  is that of IEA PAMS and CAT. Sample: the 477 country--sector cells of the 18 members of the
  shipped hand-picked list scored in both releases. Sources: \texttt{strategic\_index} of the two
  \texttt{index\_outputs.rds}; \texttt{scripts/paper/fig\_reconcile.R}.}
  \label{fig:reconcile_si_terms}
\end{figure}

\paragraph{The allied picks.} The pick that feeds the allied network is each member's top
upstream or midstream sector, Geothermal and Green Hydrogen excluded by rule of the selection
(Table~\ref{tab:reconcile_allied}). The shipped release made it over a hand-picked list of 23
economies; the final coalition is derived (Section~\ref{sec:coalition}) and has 57 members, so
the comparison is confined to the 21 members the two lists share. Kenya and the United Arab
Emirates are not in the final coalition, and Denmark, France, Italy and South Korea had no
shipped pick because they lacked a Policy index (item P3). Of the 17 members with a pick in both
releases, 16 change it; only Chile keeps Nuclear upstream. The causes separate cleanly. Six
changes (Australia, Brazil, Canada, India, Mexico and the United States) are reorderings driven
by the security-need term: all six picked Nuclear upstream, where the shipped $\ES$ was zero for
every country, and the Red Book data now record uranium for all six, which lowers their need.
Four (Argentina, Indonesia, South Africa and Viet Nam) are reorderings driven by the policy term.
Four (Germany, Japan, Morocco and Norway) are barred by the endowment screen, because the shipped
pick is an upstream stage for which the member has no recorded reserves or production. Two (Saudi
Arabia and the United Kingdom) are no longer scored, because the $\EO$ cell of the shipped pick
is not rankable. On these 21 members the picks move from Nuclear upstream (8 shipped picks) and
Batteries upstream (5) to Batteries midstream (13 final picks), with Nuclear upstream (3),
Batteries upstream and Electric Grid midstream (2 each) and Nuclear midstream (1). Over the whole
coalition every one of the 57 members has a pick: Batteries midstream 27, Nuclear midstream 9,
Nuclear upstream 8, Batteries upstream and Electric Grid midstream 4 each, Solar midstream 3 and
Wind upstream 2. Without the endowment screen Nuclear upstream would be the pick of 17 members.

\begin{table}[htbp]
  \centering
  \begin{threeparttable}
  \caption{Allied pick per member of the shipped list: shipped and final release}
  \label{tab:reconcile_allied}
  \footnotesize
  \begin{tabularx}{\textwidth}{@{}l >{\raggedright\arraybackslash}p{1.15in} c >{\raggedright\arraybackslash}p{1.15in} c >{\raggedright\arraybackslash}X@{}}
    \toprule
    & \multicolumn{2}{c}{Shipped release} & \multicolumn{2}{c}{Final release} & \\
    \cmidrule(lr){2-3}\cmidrule(lr){4-5}
    Member & Pick & $\SI$ & Pick & $\SI$ & Status or cause of change \\
    \midrule
    Argentina & Electric Vehicles, Midstream & 0.624 & Batteries, Upstream & 0.551 & reordered by the policy term \\
    Australia & Nuclear, Upstream & 0.656 & Batteries, Midstream & 0.610 & reordered by the security-need term \\
    Brazil & Nuclear, Upstream & 0.613 & Batteries, Midstream & 0.606 & reordered by the security-need term \\
    Canada & Nuclear, Upstream & 0.638 & Electric Grid, Midstream & 0.615 & reordered by the security-need term \\
    Chile & Nuclear, Upstream & 0.626 & Nuclear, Upstream & 0.645 & unchanged \\
    Denmark & -- & -- & Batteries, Midstream & 0.636 & new \\
    France & -- & -- & Batteries, Midstream & 0.642 & new \\
    Germany & Batteries, Upstream & 0.712 & Batteries, Midstream & 0.623 & shipped pick screened: no resource base \\
    India & Nuclear, Upstream & 0.643 & Batteries, Midstream & 0.628 & reordered by the security-need term \\
    Indonesia & Electric Vehicles, Midstream & 0.642 & Nuclear, Midstream & 0.568 & reordered by the policy term \\
    Italy & -- & -- & Batteries, Midstream & 0.639 & new \\
    Japan & Batteries, Upstream & 0.657 & Nuclear, Upstream & 0.620 & shipped pick screened: no resource base \\
    Kenya & Batteries, Upstream & 0.701 & \emph{not a member} & -- & not in final coalition \\
    Mexico & Nuclear, Upstream & 0.624 & Batteries, Midstream & 0.564 & reordered by the security-need term \\
    Morocco & Solar, Upstream & 0.683 & Electric Grid, Midstream & 0.609 & shipped pick screened: no resource base \\
    Norway & Nuclear, Upstream & 0.729 & Batteries, Upstream & 0.618 & shipped pick screened: no resource base \\
    Saudi Arabia & Batteries, Upstream & 0.688 & Nuclear, Upstream & 0.567 & shipped pick not scored: EO not rankable \\
    South Africa & Electric Grid, Midstream & 0.679 & Batteries, Midstream & 0.620 & reordered by the policy term \\
    South Korea & -- & -- & Batteries, Midstream & 0.635 & new \\
    United Arab Emirates & -- & -- & \emph{not a member} & -- & not in final coalition \\
    United Kingdom & Batteries, Upstream & 0.690 & Batteries, Midstream & 0.617 & shipped pick not scored: EO not rankable \\
    United States & Nuclear, Upstream & 0.613 & Batteries, Midstream & 0.632 & reordered by the security-need term \\
    Viet Nam & Batteries, Upstream & 0.698 & Batteries, Midstream & 0.606 & reordered by the policy term \\
    \bottomrule
  \end{tabularx}
  \begin{tablenotes}[flushleft]\footnotesize
    \item \emph{Notes:} The pick is the member's Upstream or Midstream sector with the highest Strategic Index; Geothermal and Green Hydrogen are excluded by rule of the selection. Shipped release: the Strategic Index of 3~September 2026 over its hand-picked list of 23 economies, of which 18 had a pick. Final release: \texttt{data/processed/charts/allied\_strategic.csv}, over the derived coalition of 57 members, with the endowment screen (an Upstream extraction stage is eligible only where the member has recorded reserves or production of its resources). The comparison is on the 21 members the two lists share; 2 economies of the shipped list are not in the final coalition. ``new'': a pick only in the final release. Cause of a change: the shipped pick is no longer scored (its exclusion), is barred by the endowment screen, or is outscored, in which case the term of the index that contributes most to the reversal of the two sectors' order is named (an exact split of the reversal into the five weighted terms). The picks of the other 36 final members are reported in Section~\ref{sec:strategic}.
  \end{tablenotes}
  \end{threeparttable}
\end{table}


\paragraph{Allied network.} The shipped network outputs are not comparable with the final ones
(Table~\ref{tab:reconcile_network}). The shipped diagnostics, the only shipped network artefact
read here, record a greedy solution in all 33 stages over the 22 effective members of the
hand-picked list, and in 14 of them a reported producer count that is a perfect square, $k^2$
with $k$ between 16 and 440 (up to 193{,}600 producers), together with a Herfindahl index of
designed shares equal to one. Both are artefacts of the development join fan-out (item N1): a
single country was selected with $k^2$ rows of share $1/k$, so the row count was $k^2$ and the
Herfindahl index summed over the rows equalled one \eqref{eq:corr_fanout}. The portfolio loop,
which runs over all stages jointly, stopped at its limit of ten passes with twelve members still a
top producer in more stages than their cap allows. The final design covers the 22 upstream and
midstream stages over 37 nodes, with the EU-27 as one node, and is solved by the MILP in every
stage with no fallback. It selects 7 to 10 producers per stage (median 9), with a designed
Herfindahl index between 0.13 and 0.21 (median 0.152, 6.6 effective producers) and no producer
above a share of 0.325, under the concave diversification term with $\lambda=0.78$; the median
observed export Herfindahl index of the nodes over the 18 stages with an HS6 basket is 0.156. The
portfolio loop converges after seven passes with no member above its cap, and the node scores
carry no policy term, so no member's policy is imputed.

\begin{table}[htbp]
  \centering
  \begin{threeparttable}
  \caption{Allied network diagnostics: shipped and final release}
  \label{tab:reconcile_network}
  \small
  \begin{tabular}{@{}l c c@{}}
    \toprule
    & Shipped release & Final release \\
    \midrule
    Network nodes & 22 & 37 \\
    Stages designed & 33 (all three stages) & 22 (Upstream, Midstream) \\
    Solver used (stages) & greedy (33) & MILP (22) \\
    Stages with $n_{\text{producers}}=k^2$ and $\HHI=1$ & 14 & 0 \\
    Reported producers per stage & 3--193{,}600 & 7--10 (median 9) \\
    $\HHI$ of designed shares & 0.333--1.000 & 0.126--0.209 (median 0.1520) \\
    Diversification weight $\lambda$ & -- & 0.78 \\
    Policy term in node scores & yes & no \\
    Portfolio loop converged & no & yes \\
    Portfolio passes & 10 (limit) & 7 \\
    Members above their portfolio cap at exit & 12 & 0 \\
    \bottomrule
  \end{tabular}
  \begin{tablenotes}[flushleft]\footnotesize
    \item \emph{Notes:} Shipped release: \texttt{allied\_network\_diagnostics.csv} written on 23~February 2026, the only shipped network artefact read here; it designed every stage of the hand-picked list's 22 effective members (the United Arab Emirates was coded ``UAE'', item N5). In 14 shipped stages a single country was selected with $k^2$ rows of share $1/k$, so the producer count was $k^2$ (with $k$ from 16 to 440) and the Herfindahl index summed over the rows was exactly one (item N1, equation~\eqref{eq:corr_fanout}). Final release: the same file of the final run of 12~October 2026, 37 nodes (the EU-27 as one node), Upstream and Midstream stages only, MILP with the concave diversification term ($\lambda=0.78$, the smallest value on the calibration grid whose median designed $\HHI$ does not exceed the observed one) and portfolio caps, and no policy term in the node scores; the median designed $\HHI$ compares with a median observed export $\HHI$ of the nodes of 0.1556 over the 18 stages with an HS6 basket. The portfolio loop runs over all stages jointly, so it has one outcome per release; the last row counts the members that, when the loop stops, are a top producer in more stages than their cap allows.
  \end{tablenotes}
  \end{threeparttable}
\end{table}


\paragraph{Partnership Strength.} The shipped release did not write the PSI to disk (item B4).
It is reconstructed here from the March 2026 top-five partner tables with the earlier code's
logic: the unit is the partner name, so Taiwan and the 18 small territories that leave the
pillar tables, none of which has a World Bank name, pool into one unnamed unit; a category score
is the mean over all reporters and all variables of the category; and a development score that
is undefined keeps its weight of 0.2 but contributes nothing. The reconstruction reproduces the
audit's figures, a mean PSI of 0.598 for partners with a development score and 0.502 for the
others. The shipped PSI averaged over every reporter, so the like-for-like comparison is with the
final world view. Over the 187 partners matched and rankable in that view the Spearman correlation
is $-0.09$ (Pearson $-0.22$); over the 191 of the headline core view, which averages the
viewpoints of the 31 core members, it is $-0.33$. The shipped ordering does not survive the
corrections, for two reasons. The denominator defect placed the partners without a development
score, mostly high-income economies, below those with one, so correcting it moves that group up
as a block. Within the two groups agreement is uneven: in the world view the Spearman correlation
is 0.69 among the 128 partners with a development score and 0.29 among the 59 without one, whose
shipped score rested on friendshore and opportunity alone, the two categories that items B1, B2
and B5 changed. The United States moves from 168th to 4th of 187 in the world view and to 1st of
191 in the core view; Germany from 135th to 5th; China from 22nd to 1st in the world view and to
44th in the core view, whose Friendshore scores rate China from the core members' viewpoints,
alignment included.

\paragraph{Using numbers from the shipped release.} Three conclusions follow for a reader who
holds numbers from the shipped release. First, within a technology--stage cell, shipped pillar
rankings are a rough guide in the $\ES$ and $\EO$ downstream, where the median within-cell
correlations are 0.90 and 0.88, and at the $\EO$ midstream stage (0.80). They are a poor guide at
the $\ES$ midstream stage (0.70) and the $\EO$ upstream stage (0.61), and they carry no
information for $\EO$ in Green Hydrogen midstream ($-0.55$). Second, comparisons across cells,
and therefore the shipped Strategic Index and its allied picks, inherit the level shifts
documented above as well as a different Policy index, and the shipped picks should be replaced by
those of Table~\ref{tab:reconcile_allied} and Section~\ref{sec:strategic}. Third, the shipped
network outputs and any PSI computed from the shipped partner tables should not be used. The
general lesson is the one the composite-indicator literature draws from uncertainty and
sensitivity analysis \citep{saisana2005uncertainty,saltelli2007composite}: the ordering a
composite produces can depend on construction choices as much as on the data. Here, however, the
large changes come from defects in the construction and from new data rather than from
defensible alternative choices, so they are corrections rather than a measure of uncertainty.
